% !TEX program = pdflatex
%% Exemple complet de la classe upbcours : Gestion de projet numérique
%% Le nom, les contacts et les données pédagogiques de l'enseignant sont fictifs.
\documentclass[12pt,twoside]{upbcours}

\upbsetup{
  titre           = {Gestion de projet numérique},
  titrecourt      = {Projet numérique},
  codes           = {UE : Management et innovation},
  soustitre       = {Licence professionnelle en systèmes d'information (Deuxième année)},
  auteurs         = {Dr Élodie KOUASSI},
  etablissement   = {Université fictive du Golfe de Guinée\\\textbf{(UFGG)}},
  faculte         = {Institut supérieur des technologies appliquées\\\textbf{(ISTA)}},
  dateproduction  = {Novembre 2026},
  datemiseajour   = {\today},
  anneeacademique = {2026 -- 2027},
  courriel        = {elodie.kouassi@ufgg-exemple.edu},
  logouniversite  = {logos/uni1.png},
  logoufr         = {logos/ecole1.png},
}

\upbsetup{
  typechapitre = {Module},
  numeroter    = oui,
  niveautdm    = 1,
  titreresume  = {À retenir},
  couleurtitres = {245B75},
  couleurfilets = {6B9DB2},
  couleuronglet = {EAF3F6},
}

\begin{document}

\frontmatter
\pagedegarde
\pagedroits
\begin{remerciements}
  \personne{Maya TCHIBOZO, relecture pédagogique}
  \personne{Cellule d'innovation et de transformation numérique de l'UFGG}
  \personne{Étudiants de la promotion pilote 2026--2027}
\end{remerciements}
\tabledesmatieres

\mainmatter
\partieupb{Syllabus}
\sylset{formation}{Licence professionnelle en systèmes d'information}
\sylset{grade}{Maître de conférences}
\sylset{domaine}{Sciences et technologies}
\sylset{annee}{2026--2027}
\sylset{semestre}{3}
\sylset{mention}{Systèmes d'information}
\sylset{specialite}{Conception et pilotage des services numériques}
\sylset{ue}{Management et innovation}
\sylset{codeec}{SIN2304}
\sylset{intituleec}{Gestion de projet numérique}
\sylset{masse}{45}
\sylset{ct}{24}
\sylset{tptd}{21}
\sylset{tpe}{30}
\sylset{credits}{4}
\sylset{effectifs}{32}
\sylset{enseignants}{Dr Élodie KOUASSI (données fictives)}
\sylset{contacts}{elodie.kouassi@ufgg-exemple.edu}
\sylset{prerequis}{Bases de conception logicielle, expression écrite, travail en équipe et
  utilisation courante d'un environnement numérique.}
\sylset{descriptif}{Ce cours accompagne les étudiants dans la préparation, la réalisation et
  l'évaluation d'un projet numérique. Il met en relation les besoins d'un commanditaire,
  l'organisation d'une équipe, la conception d'un produit utile et le suivi de sa réalisation.}
\sylset{objectifgeneral}{Conduire un projet numérique de manière structurée, collaborative,
  responsable et orientée vers la valeur produite pour ses utilisateurs.}
\sylset{objectifsspecifiques}{%
  \textbf{Savoir}\par\nobreak
  \begin{itemize}
    \item distinguer un problème, un besoin, une solution et un livrable ;
    \item décrire les rôles, les contraintes et les parties prenantes d'un projet ;
    \item comparer les approches séquentielle, agile et hybride.
  \end{itemize}
  \textbf{Savoir-faire}\par\nobreak
  \begin{itemize}
    \item cadrer un projet et formuler des objectifs vérifiables ;
    \item construire une feuille de route, un tableau de risques et un budget simplifié ;
    \item animer une réunion, documenter une décision et présenter un prototype.
  \end{itemize}
  \textbf{Savoir-être}\par\nobreak
  \begin{itemize}
    \item écouter les utilisateurs et argumenter avec respect ;
    \item rendre compte de façon fiable et prendre soin des engagements collectifs.
  \end{itemize}}
\sylset{contenu}{%
  \begin{description}[leftmargin=0pt,labelindent=0pt,style=unboxed,font=\bfseries,itemsep=3pt]
    \item[Module 1 : Comprendre le projet.] Vocabulaire, valeur, parties prenantes et critères de réussite.
    \item[Module 2 : Cadrer et planifier.] Note de cadrage, périmètre, découpage et calendrier.
    \item[Module 3 : Organiser l'équipe.] Rôles, communication, décisions et collaboration.
    \item[Module 4 : Concevoir pour les usages.] Recherche utilisateur, parcours, prototype et accessibilité.
    \item[Module 5 : Piloter l'incertitude.] Risques, qualité, budget, indicateurs et arbitrages.
    \item[Module 6 : Livrer et apprendre.] Démonstration, déploiement, bilan et amélioration continue.
  \end{description}}
\sylset{methodes}{Cours dialogué, études de cas, ateliers en petits groupes, simulations de
  réunions, revue par les pairs et projet fil rouge.}
\sylset{lieu}{Campus numérique de l'UFGG et salle de projet 204}
\sylset{materiel}{Ordinateur portable, tableau de suivi partagé, feuilles A3, cartes de persona,
  vidéoprojecteur et accès à un outil de prototypage.}
\sylset{references}{%
  \small
  \begin{enumerate}[label={[\arabic*]},leftmargin=9mm,itemsep=1.5pt]
    \item A. Osterwalder et Y. Pigneur. \emph{Business Model Nouvelle Génération}. Pearson, 2011.
    \item K. Schwaber et J. Sutherland. \emph{The Scrum Guide}. Édition pédagogique, 2020.
    \item D. C. Cagan. \emph{Inspired}. Wiley, 2017.
    \item E. Ries. \emph{Lean Startup}. Crown Business, 2011.
    \item J. Nielsen. \emph{Usability Engineering}. Morgan Kaufmann, 1994.
    \item ISO. \emph{ISO 21502 : Management de projet, lignes directrices}. 2020.
  \end{enumerate}}
\sylset{evalinter}{x}
\sylset{evalfinale}{x}
\sylset{pcinter}{40\%}
\sylset{pcfinal}{60\%}
\syllabusupb
\explicationicones

\chapitreupb{Comprendre le projet numérique}
\introduction
\begin{survol}
  Un projet n'est pas seulement une liste de tâches : c'est une transformation temporaire
  qui doit produire une valeur reconnaissable. Ce module donne le vocabulaire commun nécessaire
  avant de choisir une méthode ou un outil.
\end{survol}
\begin{objectifs}
  \item distinguer finalité, objectif, résultat et livrable ;
  \item identifier les attentes contradictoires d'un projet réel ;
  \item formuler une première hypothèse de valeur.
\end{objectifs}
\begin{concepts}
  \notion{Projet}{Démarche temporaire visant un résultat singulier dans un contexte donné.}
  \notion{Valeur}{Bénéfice attendu par une personne ou une organisation, au regard d'un coût.}
  \notion{Partie prenante}{Personne ou groupe qui influence le projet ou subit ses effets.}
\end{concepts}
\sujet{De l'idée au problème}
Une idée de service ne constitue pas encore un projet. Pour devenir actionnable, elle doit être
reliée à une situation observable, à des personnes concernées et à une amélioration attendue.
La question utile n'est donc pas « quelle application voulons-nous développer ? », mais
« quelle difficulté voulons-nous réduire, pour qui, et comment le saurons-nous ? »
\begin{definition}[Problème de projet]
  Un problème est un écart documenté entre une situation actuelle et une situation souhaitée,
  dont la réduction est possible dans les contraintes du projet.
\end{definition}
\begin{icone}{astuce}
  Remplacez les verbes vagues comme « moderniser » par un comportement observable :
  « permettre à un usager de trouver une information en moins de trois minutes ».
\end{icone}

\sujet{Les dimensions de la réussite}
La réussite ne se résume pas au respect d'une date. Un projet numérique équilibré tient compte de
la valeur d'usage, de la qualité, du coût, du délai et de l'appropriation. Ces dimensions peuvent
entrer en tension : ajouter une fonctionnalité peut augmenter la valeur perçue tout en retardant
la mise à disposition du service.
\begin{theoreme}[des contraintes liées]
  Toute décision de projet modifie au moins une combinaison de périmètre, délai, ressources,
  qualité et risques. Une décision fiable explicite le compromis accepté.
\end{theoreme}
\begin{exemple}
  Pour un portail d'inscription, l'équipe peut livrer d'abord la consultation des pièces
  nécessaires, puis reporter la prise de rendez-vous à une version ultérieure.
\end{exemple}
\sousujet{Carte des parties prenantes}
Classez les acteurs selon leur pouvoir d'influence et leur niveau d'intérêt. Le commanditaire,
l'équipe, les utilisateurs finaux et les services support n'ont pas les mêmes critères de réussite.
\sujet{Atelier : la fiche valeur}
\begin{icone}{activite}
  En groupes de quatre, choisissez un service public fictif. Décrivez un problème en une phrase,
  trois utilisateurs concernés, un bénéfice mesurable et une contrainte non négociable.
\end{icone}
\begin{enumerate}
  \item Décrire la situation actuelle sans proposer de solution.
  \item Recueillir deux faits qui confirment le problème.
  \item Écrire une hypothèse : « Nous pensons que \ldots parce que \ldots ».
  \item Proposer un indicateur simple qui permettrait de l'infirmer.
\end{enumerate}
\begin{resumechapitre}
  Un projet numérique part d'un problème situé et produit une valeur vérifiable. Les contraintes
  doivent être rendues visibles afin que les arbitrages futurs soient compréhensibles.
\end{resumechapitre}
\begin{evalformative}
  Transformez « créer une plateforme moderne » en objectif observable pour un service de bibliothèque.
\end{evalformative}
\begin{autoevaluation}
  \begin{enumerate}
    \item Quelle différence faites-vous entre un livrable et un bénéfice ?
    \item Citez quatre parties prenantes d'un projet numérique.
    \item Pourquoi une hypothèse de valeur doit-elle pouvoir être infirmée ?
  \end{enumerate}
\end{autoevaluation}

\chapitreupb{Cadrer et planifier}
\introduction
\begin{survol}
  Le cadrage transforme une intention en contrat de travail lisible. Il fixe le pourquoi, le pour
  qui, le quoi et le comment sans prétendre supprimer toute incertitude.
\end{survol}
\begin{objectifs}
  \item rédiger une note de cadrage courte ;
  \item définir un périmètre inclus et exclu ;
  \item construire une première feuille de route.
\end{objectifs}
\sujet{La note de cadrage}
Une note de cadrage tient sur quelques pages et sert de référence commune. Elle présente le contexte,
le problème, les résultats attendus, les exclusions, les hypothèses, les acteurs et le mode de décision.
\begin{definition}[Périmètre]
  Le périmètre est l'ensemble des résultats que le projet s'engage à produire, accompagné de ses
  limites explicites.
\end{definition}
\begin{icone}{miseengarde}
  Une exclusion écrite n'est pas un refus définitif : c'est une protection contre l'ambiguïté
  pendant la période où l'équipe doit tenir ses engagements.
\end{icone}
\sujet{Découper le travail}
Le découpage en résultats puis en lots rend le travail discut able. Chaque lot doit avoir un
responsable, une condition de fin et une dépendance connue. Évitez les tâches trop générales
comme « faire le site » ; préférez « valider le formulaire de demande avec trois utilisateurs ».
\begin{exemple}
  Pour un catalogue en ligne : recueillir les critères, dessiner le parcours, prototyper la
  recherche, tester le prototype, puis préparer la mise en ligne.
\end{exemple}
\begin{icone}{note}
  La planification est une hypothèse révisable, pas une promesse de prédire exactement l'avenir.
\end{icone}
\sujet{Calendrier et jalons}
Un jalon marque une décision ou un résultat vérifiable. Il ne correspond pas nécessairement à une
date de fin de tâche. La feuille de route peut présenter trois horizons : maintenant, ensuite,
et plus tard. Cette vue évite de détailler artificiellement des travaux encore mal compris.
\begin{table}[ht]
  \centering
  \begin{tabular}{p{0.25\linewidth}p{0.28\linewidth}p{0.32\linewidth}}
    \toprule
    Horizon & Résultat & Question de décision\\
    \midrule
    Maintenant & Parcours prioritaire prototypé & Le besoin est-il compris ?\\
    Ensuite & Version testable & Les utilisateurs réussissent-ils la tâche ?\\
    Plus tard & Service déployé & Le service est-il soutenable ?\\
    \bottomrule
  \end{tabular}
  \caption{Exemple de feuille de route à trois horizons.}
\end{table}
\begin{resumechapitre}
  Le cadrage rend les engagements discutables et les limites visibles. Un bon plan relie des
  résultats à des décisions, plutôt qu'une simple succession de tâches.
\end{resumechapitre}
\begin{evalformative}
  Rédigez le périmètre inclus et exclu d'un outil fictif de réservation de salles.
\end{evalformative}
\begin{autoevaluation}
  \begin{enumerate}
    \item Quelles rubriques doivent figurer dans une note de cadrage ?
    \item Comment reconnaître une tâche trop vague ?
    \item Quelle est l'utilité d'un jalon ?
  \end{enumerate}
\end{autoevaluation}

\chapitreupb{Organiser l'équipe et la communication}
\introduction
\begin{survol}
  Les projets avancent par conversations, décisions et apprentissages. Ce module propose des
  repères pour distribuer les responsabilités sans isoler les compétences.
\end{survol}
\begin{objectifs}
  \item clarifier les responsabilités ;
  \item choisir un canal et un rythme de communication ;
  \item conduire une réunion courte et utile.
\end{objectifs}
\sujet{Rôles et responsabilités}
Une équipe efficace sait qui prépare, qui décide, qui contribue et qui doit être informé. Une même
personne peut tenir plusieurs rôles dans une petite équipe, mais la responsabilité d'une décision
doit rester explicite.
\begin{definition}[Responsabilité]
  Une responsabilité est un engagement observable envers un résultat, avec l'autorité minimale
  nécessaire pour agir ou demander un arbitrage.
\end{definition}
\begin{icone}{saviezvous}
  Dans une petite équipe, le rôle de facilitateur peut tourner chaque semaine afin de distribuer
  la charge invisible de préparation et de suivi.
\end{icone}
\sujet{La réunion comme outil de production}
Une réunion utile possède un objectif, une durée, des participants justifiés et une trace. Commencez
par le résultat attendu, distinguez les informations des décisions, puis terminez par les actions
attribuées. Le compte rendu doit pouvoir être compris par une personne absente.
\begin{enumerate}
  \item Ouverture : objectif et contraintes de temps.
  \item Exploration : faits, questions et options.
  \item Décision : choix, responsable et date de réexamen.
  \item Clôture : synthèse et diffusion de la trace.
\end{enumerate}
\sujet{Gérer les désaccords}
Un désaccord explicite est souvent plus sain qu'un accord silencieux. Reformulez d'abord le point
de vue de l'autre, séparez les faits des préférences et proposez une expérience limitée quand la
discussion ne permet pas de trancher.
\begin{icone}{jeu}
  Simulation : une équipe doit choisir entre une livraison rapide et une fonctionnalité complète.
  Un étudiant joue le commanditaire, deux l'équipe et un autre la personne utilisatrice.
\end{icone}
\begin{resumechapitre}
  L'organisation d'une équipe repose sur des responsabilités visibles, des décisions tracées et
  des échanges adaptés à l'urgence. La communication est une partie du travail, pas une décoration.
\end{resumechapitre}
\begin{evalformative}
  Préparez l'ordre du jour d'une réunion de quinze minutes destinée à débloquer une tâche.
\end{evalformative}
\begin{autoevaluation}
  \begin{enumerate}
    \item Pourquoi attribuer un responsable à une action ?
    \item Quelle différence entre information et décision ?
    \item Comment rendre un désaccord productif ?
  \end{enumerate}
\end{autoevaluation}

\chapitreupb{Concevoir pour les usages}
\introduction
\begin{survol}
  Concevoir pour les usages consiste à comprendre des situations réelles avant de figer une
  solution. Les prototypes rendent les idées discutables à faible coût.
\end{survol}
\begin{objectifs}
  \item préparer un entretien utilisateur ;
  \item représenter un parcours et une priorité ;
  \item construire un prototype compréhensible et accessible.
\end{objectifs}
\begin{concepts}
  \notion{Persona}{Portrait synthétique d'un groupe d'utilisateurs partageant une situation.}
  \notion{Parcours}{Suite d'étapes et d'émotions qui relie un besoin à un résultat.}
  \notion{Prototype}{Représentation volontairement incomplète destinée à apprendre.}
\end{concepts}
\sujet{Observer avant de concevoir}
Un entretien ne cherche pas à faire valider une solution déjà choisie. Il recueille des exemples
concrets : dernière occurrence, contournements, vocabulaire, contraintes et critères de confiance.
Les questions ouvertes sont préférables aux questions qui suggèrent la réponse.
\begin{icone}{lecture}
  Lisez une page de service public et repérez les mots que l'institution emploie, puis ceux
  qu'une personne novice aurait besoin de comprendre.
\end{icone}
\sujet{Du parcours au prototype}
Le parcours décrit le contexte avant l'écran. Pour chaque étape, notez l'intention de la personne,
son action, l'information dont elle a besoin et le signal qui confirme la réussite. Le prototype
peut ensuite représenter uniquement les moments qui portent le plus d'incertitude.
\begin{exemple}
  Une demande de document peut être représentée par quatre écrans : comprendre les pièces,
  vérifier l'éligibilité, déposer le dossier et obtenir une confirmation traçable.
\end{exemple}
\sujet{Accessibilité et critères d'acceptation}
L'accessibilité s'anticipe dès le cadrage : contraste, navigation au clavier, langage clair,
taille des zones interactives et alternatives aux contenus visuels. Un critère d'acceptation décrit
un comportement vérifiable plutôt qu'une impression générale de qualité.
\begin{theoreme}[du prototype honnête]
  Plus une décision est coûteuse à modifier, plus elle doit être testée sous une forme simple
  avant son implémentation complète.
\end{theoreme}
\begin{resumechapitre}
  L'écoute produit des faits, le parcours donne du contexte et le prototype permet d'apprendre
  avant d'investir. L'accessibilité et les critères de réussite font partie de la conception.
\end{resumechapitre}
\begin{evalformative}
  Écrivez trois questions d'entretien qui explorent une inscription récente sans proposer de solution.
\end{evalformative}
\begin{autoevaluation}
  \begin{enumerate}
    \item Qu'est-ce qu'un prototype doit permettre d'apprendre ?
    \item Pourquoi partir d'un exemple vécu pendant un entretien ?
    \item Donnez deux critères d'acceptation vérifiables.
  \end{enumerate}
\end{autoevaluation}

\chapitreupb{Piloter l'incertitude}
\introduction
\begin{survol}
  Piloter ne signifie pas tout prévoir. Il s'agit de rendre les incertitudes visibles, de choisir
  les vérifications les plus utiles et de décider à partir d'informations suffisamment fiables.
\end{survol}
\begin{objectifs}
  \item décrire un risque par sa cause et son effet ;
  \item choisir une réponse proportionnée ;
  \item suivre quelques indicateurs sans suradministrer le projet.
\end{objectifs}
\sujet{Le registre des risques}
Un risque est un événement incertain dont la survenue peut modifier les objectifs. Décrivez sa cause,
son événement, son effet, sa probabilité et sa gravité. Ajoutez un propriétaire et une prochaine
action de réduction : un registre sans action n'est qu'une liste d'inquiétudes.
\begin{definition}[Risque]
  Un risque est une possibilité future, différente d'un problème déjà avéré, qui mérite une décision
  de prévention, de réduction, de transfert ou d'acceptation.
\end{definition}
\begin{icone}{miseengarde}
  Ne transformez pas une échelle de couleur en vérité mathématique : discutez ce que chaque niveau
  signifie dans le contexte précis du projet.
\end{icone}
\sujet{Qualité et indicateurs}
La qualité se vérifie à partir de critères définis avant la démonstration. Un indicateur doit avoir
un nom, une source, une fréquence, une cible et une personne responsable de son interprétation.
Comptez moins d'indicateurs, mais utilisez-les pour agir.
\begin{table}[ht]
  \centering
  \begin{tabular}{p{0.3\linewidth}p{0.25\linewidth}p{0.3\linewidth}}
    \toprule
    Indicateur & Source & Décision possible\\
    \midrule
    Taux de réussite d'une tâche & Test observé & Revoir le parcours\\
    Délai de réponse & Journal du service & Simplifier le traitement\\
    Défauts bloquants & Revue de version & Retarder la livraison\\
    \bottomrule
  \end{tabular}
  \caption{Indicateurs orientés vers l'action.}
\end{table}
\sujet{Arbitrer avec transparence}
Un arbitrage est une décision qui accepte une perte relative pour protéger une priorité. Écrivez
le contexte, les options étudiées, le choix, la raison, les effets attendus et la date de réexamen.
Cette trace évite de rouvrir sans cesse la même discussion.
\begin{icone}{devoir}
  À partir d'un registre de cinq risques, préparez une note d'arbitrage d'une page et défendez-la
  devant un autre groupe.
\end{icone}
\begin{resumechapitre}
  Un risque devient pilotable quand il est décrit, attribué et relié à une action. Les indicateurs
  n'ont de valeur que s'ils éclairent une décision et si leurs limites sont connues.
\end{resumechapitre}
\begin{evalformative}
  Formulez un risque lié à une dépendance externe et proposez deux réponses de réduction.
\end{evalformative}
\begin{autoevaluation}
  \begin{enumerate}
    \item Comment distinguer un risque d'un problème ?
    \item Quelles informations rendent un indicateur actionnable ?
    \item Que faut-il conserver dans une décision d'arbitrage ?
  \end{enumerate}
\end{autoevaluation}

\chapitreupb{Livrer, évaluer et apprendre}
\introduction
\begin{survol}
  La livraison est une transition vers l'usage, pas la fin de la responsabilité. Une démonstration,
  une mise en service accompagnée et un bilan honnête permettent de transformer l'expérience en progrès.
\end{survol}
\begin{objectifs}
  \item préparer une démonstration centrée sur la valeur ;
  \item organiser une mise en service progressive ;
  \item rédiger un bilan qui distingue faits, interprétations et décisions.
\end{objectifs}
\sujet{Démontrer la valeur}
Une démonstration raconte une situation avant de montrer une fonctionnalité. Elle présente le
problème, le parcours, le résultat observable et les limites connues. Une bonne démonstration
permet au public de poser des questions précises, même lorsque le produit est incomplet.
\begin{icone}{guide}
  Préparez un scénario nominal, un scénario d'erreur et une réponse à la question :
  « Qu'avez-vous appris qui modifie la suite du projet ? »
\end{icone}
\sujet{Déployer avec prudence}
Une mise en service progressive réduit le coût d'une erreur. Définissez qui peut accéder au service,
comment revenir en arrière, quels signaux surveiller et qui répond aux utilisateurs. Le transfert
de compétences et la documentation sont des livrables à part entière.
\begin{exemple}
  Une équipe ouvre son formulaire à dix utilisateurs volontaires pendant deux jours, observe les
  demandes d'aide, corrige les obstacles puis élargit l'accès.
\end{exemple}
\sujet{Bilan et amélioration continue}
Le bilan ne cherche pas un coupable. Il compare les hypothèses initiales aux observations, identifie
les décisions utiles et choisit quelques changements concrets. Une rétrospective réussie produit
des engagements modestes, attribués et datés.
\begin{icone}{evaluation}
  Évaluation du projet fil rouge : dossier de cadrage (20\%), prototype testé (25\%), présentation
  finale (25\%), journal d'équipe et rétrospective (30\%).
\end{icone}
\begin{resumechapitre}
  Livrer consiste à rendre une valeur utilisable et soutenable. La démonstration, le déploiement
  progressif et le bilan transforment un résultat ponctuel en capacité d'apprendre.
\end{resumechapitre}
\begin{evalformative}
  Préparez une liste de vérification pour la mise en service d'un service fictif.
\end{evalformative}
\begin{autoevaluation}
  \begin{enumerate}
    \item Qu'est-ce qu'une démonstration centrée sur la valeur ?
    \item Pourquoi prévoir un retour en arrière ?
    \item Comment distinguer un fait d'une interprétation dans un bilan ?
  \end{enumerate}
\end{autoevaluation}

\partieupb{Ressources et travaux dirigés}
\sujet{Projet fil rouge : Guichet Clair}
Les équipes conçoivent un parcours fictif permettant à des étudiants de trouver une aide
administrative adaptée. Le projet doit rester limité à l'information et à l'orientation :
aucune donnée personnelle réelle ne doit être collectée.
\begin{objectifs}
  \item produire une note de cadrage de deux pages ;
  \item réaliser un prototype de cinq écrans maximum ;
  \item présenter trois enseignements tirés de tests utilisateurs fictifs.
\end{objectifs}
\sujet{Barème de la présentation finale}
\begin{tabular}{p{0.55\linewidth}r}
  \toprule
  Critère & Points\\
  \midrule
  Compréhension du problème et des utilisateurs & 5\\
  Cohérence du périmètre et des décisions & 5\\
  Qualité du prototype et accessibilité & 5\\
  Argumentation, trace et retour d'expérience & 5\\
  \bottomrule
\end{tabular}

\partieupb{Annexes}
\fichetp{1}
\fichetd[
  titre       = {Cadrer le projet Guichet Clair},
  duree       = {2 h},
  objectif    = {Formuler un problème, un périmètre et trois critères de réussite.},
  equipe      = {Groupes de 4 étudiants.},
  individuel  = {Chaque étudiant rédige une hypothèse de valeur.},
  materiel    = {Fiche A3, marqueurs, tableau de suivi partagé.},
  consignes   = {\begin{enumerate}\item Décrire la situation actuelle.\item Cartographier les acteurs.
    \item Rédiger la note de cadrage.\item Faire relire le périmètre par un autre groupe.\end{enumerate}},
  restitution = {Note de cadrage et présentation de cinq minutes.},
  role        = {Questionner les hypothèses et aider à préciser les critères.},
  evaluation  = {Barème sur 20 : problème (5), périmètre (5), critères (5), coopération (5).}
]{2}
\fichetd[
  titre       = {Tester un prototype},
  duree       = {3 h},
  objectif    = {Observer un parcours et proposer des améliorations argumentées.},
  equipe      = {Binômes : une personne observe, une personne anime.},
  individuel  = {Journal réflexif d'une page.},
  materiel    = {Prototype papier, grille d'observation, chronomètre.},
  consignes   = {\begin{enumerate}\item Préparer deux scénarios.\item Observer sans aider trop tôt.
    \item Noter les faits.\item Prioriser trois améliorations.\end{enumerate}},
  restitution = {Carte des observations et prototype révisé.},
  role        = {Garantir un cadre d'observation respectueux.},
  evaluation  = {Barème sur 20 : observation (8), analyse (6), priorisation (6).}
]{3}

\end{document}
